\documentclass[aspectratio=169]{beamer}

% Tema UFS — Universidade Federal de Sergipe
\usetheme{UFS}

% Pacotes base
\usepackage[utf8]{inputenc}
\usepackage[portuguese]{babel}
\usepackage{amsmath,amssymb}
\usepackage{graphicx}
\usepackage{booktabs}
\usepackage{xcolor}

% TikZ e bibliotecas
\usepackage{tikz}
\usetikzlibrary{shapes.geometric, arrows.meta, positioning, matrix, fit,
  backgrounds, decorations.pathreplacing, shadows, patterns}

% Listagens — paleta VS Code Light+
\usepackage{listings}
\definecolor{bgcolor}{HTML}{FFFFFF}
\definecolor{linecolor}{HTML}{DDDDDD}
\definecolor{numcolor}{HTML}{999999}
\definecolor{kwcolor}{HTML}{0000FF}
\definecolor{strcolor}{HTML}{A31515}
\definecolor{cmtcolor}{HTML}{008000}

\lstdefinestyle{ufsStyle}{
  backgroundcolor=\color{bgcolor},
  basicstyle=\ttfamily\footnotesize\color{black},
  keywordstyle=\color{kwcolor}\bfseries,
  stringstyle=\color{strcolor},
  commentstyle=\color{cmtcolor}\itshape,
  numberstyle=\tiny\color{numcolor},
  numbers=left, numbersep=12pt,
  xleftmargin=20pt, framexleftmargin=20pt,
  breaklines=true, tabsize=4,
  keepspaces=true, showspaces=false,
  showstringspaces=false, showtabs=false,
  frame=single, rulecolor=\color{linecolor},
  framesep=4pt, captionpos=b, upquote=true,
  columns=fullflexible,
}
\lstset{style=ufsStyle, language=C}

% --- Metadados ---
\title{Registros (\texttt{struct})}
\subtitle{Declaração, acesso, registros aninhados e o layout dos dados na memória}
\author{Prof.\ André Yoshiaki Kashiwabara}
\institute{DCOMP -- Universidade Federal de Sergipe\\
\texttt{andreyoshiaki@dcomp.ufs.br}}
\date{}

\begin{document}

\begin{frame}[plain]
  \titlepage
\end{frame}

% =====================================================================
\section{Um programa para começar}

\begin{frame}{Por que precisamos de registros?}
  \begin{columns}[T]
    \begin{column}{0.48\textwidth}
      \textbf{Problema:}\\[2mm]
      Um aluno tem nome, matrícula, data de nascimento e média. Com o que
      sabemos até aqui, guardamos isso em \textbf{vetores paralelos}: um para
      nomes, outro para matrículas, outro para médias.
      \\[2mm]
      Nada garante que \texttt{nome[3]} e \texttt{media[3]} sejam da
      \emph{mesma} pessoa. Ordenar por média exige reordenar todos os vetores
      em sincronia.
    \end{column}
    \begin{column}{0.48\textwidth}
      \textbf{Solução:}\\[2mm]
      Um \textbf{registro} agrupa campos de tipos diferentes sob um único
      nome. O aluno passa a ser \emph{um} valor: cabe em uma variável, entra
      em um vetor, vira parâmetro e vira retorno de função.
      \\[2mm]
      ``Os dados de um aluno'' deixam de ser convenção na cabeça do
      programador e viram um \textbf{tipo} que o compilador conhece.
    \end{column}
  \end{columns}
\end{frame}

\begin{frame}{Vetores paralelos $\times$ vetor de registros}
  \begin{center}
  \resizebox{0.95\textwidth}{!}{%
  \begin{tikzpicture}[
    cel/.style={rectangle, draw=UFSGray, fill=white, minimum width=1.5cm,
      minimum height=0.6cm, font=\scriptsize},
    rot/.style={rectangle, draw=UFSBlue, fill=UFSBlue!10, minimum width=1.5cm,
      minimum height=0.6cm, font=\scriptsize},
    rot2/.style={rectangle, draw=UFSBlue, fill=UFSYellow!25, minimum width=1.5cm,
      minimum height=0.6cm, font=\scriptsize},
    rot3/.style={rectangle, draw=UFSBlue, fill=UFSRed!12, minimum width=1.5cm,
      minimum height=0.6cm, font=\scriptsize},
    rotulo/.style={font=\scriptsize\ttfamily, anchor=east}
  ]
    % --- esquerda: vetores paralelos
    \node[font=\small\bfseries] at (2.3,1.2) {Vetores paralelos};
    \node[rotulo] at (-0.2,0.4)  {nome};
    \node[rotulo] at (-0.2,-0.3) {matricula};
    \node[rotulo] at (-0.2,-1.0) {media};
    \foreach \x/\nm/\mt/\md in {0.75/Ana/202600123/8.75,
                                2.25/Bia/202600124/6.50,
                                3.75/Caio/202600125/9.10} {
      \node[cel] at (\x,0.4)  {\nm};
      \node[cel] at (\x,-0.3) {\mt};
      \node[cel] at (\x,-1.0) {\md};
    }
    \node[font=\scriptsize, text=UFSRed, align=center] at (2.3,-1.9)
      {sincronia mantida ``na mão'':\\trocar um vetor exige trocar os outros};

    % --- direita: vetor de registros
    \node[font=\small\bfseries] at (11,1.2) {Vetor de registros};
    \node[rotulo] at (8.0,-0.3) {turma};
    \foreach \i/\x in {0/9.0, 1/11.0, 2/13.0} {
      \node[rot]  at (\x,0.4)  {nome};
      \node[rot2] at (\x,-0.3) {matricula};
      \node[rot3] at (\x,-1.0) {media};
    }
    \foreach \x in {9.0, 11.0, 13.0} {
      \draw[UFSBlue, thick, rounded corners]
        (\x-0.82,-1.38) rectangle (\x+0.82,0.78);
    }
    \node[font=\scriptsize, text=UFSBlue, align=center] at (11,-1.9)
      {cada aluno é \emph{um} valor:\\\texttt{turma[0]}, \texttt{turma[1]}, \texttt{turma[2]}};
  \end{tikzpicture}}
  \end{center}
\end{frame}

\begin{frame}[fragile]{Um programa para começar: a ficha do aluno}
\begin{columns}[T]
\begin{column}{0.48\textwidth}
\begin{lstlisting}[basicstyle=\ttfamily\scriptsize]
#include <stdio.h>
#include <string.h>

struct Data {
    int dia;
    int mes;
    int ano;
};
struct Aluno {
    char        nome[20];
    int         matricula;
    struct Data nascimento;
    double      media;
};
\end{lstlisting}
\end{column}
\begin{column}{0.48\textwidth}
\begin{lstlisting}[basicstyle=\ttfamily\scriptsize, firstnumber=15]
int main(void) {
    struct Aluno a;
    strcpy(a.nome, "Ana Souza");
    a.matricula = 202600123;
    a.nascimento.dia = 7;
    a.nascimento.mes = 3;
    a.nascimento.ano = 2005;
    a.media = 8.75;
    printf("%s (%d)\n", a.nome, a.matricula);
    printf("Media: %.2f\n", a.media);
    printf("sizeof(Aluno) = %zu\n",
           sizeof(struct Aluno));
    return 0;
}
\end{lstlisting}
\end{column}
\end{columns}
\end{frame}

\begin{frame}{O que este programa faz?}
  \begin{block}{Em duplas, anotem a previsão (3--4 min) --- sem computador!}
    \begin{enumerate}
      \item Quais são exatamente as três linhas impressas na tela?
      \item Quantos \emph{bytes} ocupa \texttt{struct Data}? E \texttt{struct Aluno}?
      \item Some o tamanho de cada campo de \texttt{struct Aluno} separadamente.
        Esse total tem que bater com \texttt{sizeof(struct Aluno)}?
    \end{enumerate}
  \end{block}
  \begin{exampleblock}{Regra do jogo}
    Escreva um número, não um ``mais ou menos''. Previsão errada é o ponto de
    partida da investigação --- e a pergunta 3 é a mais interessante da aula.
  \end{exampleblock}
\end{frame}

\begin{frame}[fragile]{Compile, execute e compare}
\begin{lstlisting}[language={}, numbers=none, basicstyle=\ttfamily\small]
$ gcc ficha_aluno.c -o ficha_aluno
$ ./ficha_aluno
Ana Souza (202600123)
Media: 8.75
sizeof(Aluno) = 48
\end{lstlisting}
  \begin{alertblock}{A soma dos campos não fecha}
    \texttt{char nome[20]} $=20$, \texttt{int matricula} $=4$,
    \texttt{struct Data} $=12$, \texttt{double media} $=8$.
    Total: \textbf{44 bytes}. Mas \texttt{sizeof} respondeu \textbf{48}.
  \end{alertblock}
  \begin{block}{A pergunta que guia o resto da aula}
    Onde estão os 4 bytes que ninguém pediu --- e quem decidiu colocá-los ali?
  \end{block}
\end{frame}

% =====================================================================
\section{Registros: declaração e acesso}

\begin{frame}{Registro: definição}
  \begin{block}{Definição}
    Um \textbf{registro} (\texttt{struct}) é um tipo composto que agrupa um
    número fixo de \textbf{campos} (ou membros), possivelmente de tipos
    diferentes, sob um único nome. Cada campo tem nome próprio e ocupa uma
    região distinta da memória do registro.
  \end{block}
  \begin{exampleblock}{Intuição}
    Um vetor é uma \emph{fila de itens iguais} acessados por índice; um
    registro é uma \emph{ficha de campos diferentes} acessados por nome.
    O vetor responde ``qual é o terceiro?''; o registro responde ``qual é a média?''.
  \end{exampleblock}
\end{frame}

\begin{frame}[fragile]{Anatomia da declaração}
\begin{lstlisting}[basicstyle=\ttfamily\small]
struct Aluno {        /* Aluno e' a TAG do registro   */
    char   nome[20];  /* campo 1                      */
    int    matricula; /* campo 2                      */
    double media;     /* campo 3                      */
};                    /* ; obrigatorio aqui!          */

struct Aluno a;       /* declara UMA variavel do tipo  */
struct Aluno turma[40];/* vetor de 40 registros        */
\end{lstlisting}
  \begin{itemize}
    \item A declaração cria um \textbf{tipo}, não uma variável: nenhuma memória
      é reservada até você escrever \texttt{struct Aluno a;}.
    \item O nome completo do tipo é \texttt{struct Aluno} --- as duas palavras.
    \item Esquecer o \texttt{;} depois do \texttt{\}} é o erro mais comum.
  \end{itemize}
\end{frame}

\begin{frame}[fragile]{Um apelido para o tipo: \texttt{typedef}}
\begin{columns}[T]
\begin{column}{0.5\textwidth}
\begin{lstlisting}[basicstyle=\ttfamily\scriptsize]
/* sem typedef */
struct Aluno {
    char nome[20];
    int  matricula;
};

struct Aluno a;
struct Aluno turma[40];
\end{lstlisting}
\end{column}
\begin{column}{0.5\textwidth}
\begin{lstlisting}[basicstyle=\ttfamily\scriptsize]
/* com typedef */
typedef struct {
    char nome[20];
    int  matricula;
} Aluno;

Aluno a;
Aluno turma[40];
\end{lstlisting}
\end{column}
\end{columns}
  \begin{block}{Qual usar?}
    As duas formas geram exatamente o mesmo tipo em memória. O \texttt{typedef}
    encurta o código; a forma com \texttt{struct} deixa explícito que o tipo é
    um registro. Nesta aula usaremos \texttt{struct Nome} para não esconder nada.
  \end{block}
\end{frame}

\begin{frame}[fragile]{Acesso aos campos: o operador ponto}
\begin{lstlisting}[basicstyle=\ttfamily\small]
struct Aluno a;

a.matricula = 202600123;      /* escrita: campo a esquerda do = */
a.media     = 8.75;
strcpy(a.nome, "Ana Souza");  /* vetor de char: use strcpy      */

printf("%d\n", a.matricula);  /* leitura                        */
a.media = a.media + 0.5;      /* campo se usa como variavel     */
\end{lstlisting}
  \begin{itemize}
    \item \texttt{a.media} \emph{é} uma variável \texttt{double}: vale em toda
      expressão onde um \texttt{double} valeria.
    \item \texttt{a.nome} é um vetor --- \texttt{a.nome = "Ana"} \textbf{não}
      compila; use \texttt{strcpy}.
  \end{itemize}
\end{frame}

\begin{frame}[fragile]{Inicialização}
\begin{lstlisting}[basicstyle=\ttfamily\scriptsize]
/* 1. por lista, na ordem de declaracao dos campos */
struct Aluno a = {"Ana Souza", 202600123, 8.75};
/* 2. campos nomeados (C99) -- ordem livre, resto zerado */
struct Aluno b = {.matricula = 202600124, .media = 7.0};
/* 3. tudo zerado */
struct Aluno c = {0};
/* 4. campo a campo, depois da declaracao */
struct Aluno d;
d.matricula = 202600125;
\end{lstlisting}
  \begin{alertblock}{Cuidado}
    Sem inicialização, os campos de \texttt{d} contêm \textbf{lixo} --- o que
    estava naquela região da pilha. Ler antes de escrever dá valores
    imprevisíveis.
  \end{alertblock}
\end{frame}

\begin{frame}[fragile]{Registro copia inteiro --- vetor não}
\begin{lstlisting}[basicstyle=\ttfamily\scriptsize]
struct Aluno a = {"Ana Souza", 202600123, 8.75};
struct Aluno b;

b = a;              /* copia OS 48 BYTES: b e a sao independentes */
b.media = 10.0;     /* nao afeta a.media                          */

int v[3] = {1,2,3}, w[3];
w = v;              /* ERRO de compilacao: vetor nao se atribui   */
\end{lstlisting}
  \begin{block}{Consequência prática}
    O registro é um \textbf{valor}: pode ser atribuído, passado por cópia e
    devolvido por \texttt{return}. Repare que o vetor \texttt{nome[20]}
    \emph{dentro} do registro também é copiado --- ``vetor não se atribui''
    vale para o vetor solto, não para o campo.
  \end{block}
\end{frame}

\begin{frame}[fragile]{Registros em funções}
\begin{lstlisting}[basicstyle=\ttfamily\tiny]
/* recebe uma COPIA: alteracoes nao saem da funcao */
void imprime(struct Aluno x) {
    printf("%-20s %d  %.2f\n", x.nome, x.matricula, x.media);
}

/* devolve um registro construido dentro da funcao */
struct Aluno cria(const char *nome, int mat, double media) {
    struct Aluno novo;
    strcpy(novo.nome, nome);
    novo.matricula = mat;
    novo.media     = media;
    return novo;               /* copia de volta para quem chamou */
}
\end{lstlisting}
  \begin{exampleblock}{Guarde esta pergunta}
    E se o registro tivesse 500 bytes e fosse passado assim dentro de um laço de
    um milhão de iterações? A resposta --- passar o \emph{endereço} --- é a
    próxima aula.
  \end{exampleblock}
\end{frame}

\begin{frame}[fragile]{Vetor de registros}
\begin{lstlisting}[basicstyle=\ttfamily\scriptsize]
struct Aluno turma[3] = {
    {"Ana Souza",   202600123, 8.75},
    {"Bia Rocha",   202600124, 6.50},
    {"Caio Lima",   202600125, 9.10}
};

double soma = 0.0;
for (int i = 0; i < 3; i++)
    soma += turma[i].media;    /* indexa o vetor, DEPOIS pega o campo */

printf("Media da turma: %.2f\n", soma / 3);
\end{lstlisting}
  \begin{itemize}
    \item \texttt{turma[i].media}: primeiro \texttt{[i]} escolhe o registro,
      depois \texttt{.media} escolhe o campo dentro dele.
    \item Ordenar a turma por média agora é trocar registros inteiros --- e a
      ficha de cada aluno viaja junta, sem risco de dessincronizar.
  \end{itemize}
\end{frame}

\begin{frame}{Recapitulando: declaração e acesso}
  \begin{block}{O que aprendemos}
    \begin{itemize}
      \item \texttt{struct Tag \{ campos \};} cria um \textbf{tipo}; a memória
        só aparece ao declarar uma variável.
      \item Campos se acessam por nome com o operador ponto e valem como
        variáveis comuns.
      \item Inicialização: por lista, por campos nomeados, com \texttt{\{0\}} ou
        campo a campo --- sem inicializar, há lixo.
      \item Registros se atribuem, se passam para funções e se retornam ---
        sempre por \textbf{cópia}.
    \end{itemize}
  \end{block}
\end{frame}

% =====================================================================
\section{Registros aninhados}

\begin{frame}{Por que aninhar registros?}
  \begin{columns}[T]
    \begin{column}{0.48\textwidth}
      \textbf{Problema:}\\[2mm]
      A data de nascimento tem três partes que só fazem sentido juntas.
      Espalhá-las como \texttt{dia\_nasc}, \texttt{mes\_nasc},
      \texttt{ano\_nasc} dentro de \texttt{struct Aluno} funciona --- até
      aparecer a data de matrícula, e depois a de colação. Aí são nove campos
      soltos e nenhuma função que receba ``uma data''.
    \end{column}
    \begin{column}{0.48\textwidth}
      \textbf{Solução:}\\[2mm]
      Declarar \texttt{struct Data} uma vez e \textbf{usá-la como tipo de
      campo}. Toda função que valida, compara ou imprime datas passa a servir
      a qualquer registro que contenha uma --- e o conceito ``data'' fica
      escrito em um lugar só.
    \end{column}
  \end{columns}
\end{frame}

\begin{frame}[fragile]{Um registro como campo de outro}
\begin{lstlisting}[basicstyle=\ttfamily\scriptsize]
struct Data {              /* precisa vir ANTES de ser usada */
    int dia, mes, ano;
};

struct Aluno {
    char        nome[20];
    int         matricula;
    struct Data nascimento;   /* campo do tipo struct Data */
    struct Data matriculado;  /* outro campo, mesmo tipo   */
    double      media;
};
\end{lstlisting}
  \begin{itemize}
    \item \texttt{struct Data} tem de estar declarada antes --- o compilador
      precisa saber o tamanho do campo.
    \item Um registro \textbf{não} pode conter a si mesmo (seria de tamanho
      infinito); só um \emph{ponteiro} para si mesmo.
  \end{itemize}
\end{frame}

\begin{frame}{Acesso encadeado: leia da esquerda para a direita}
  \begin{center}
  \resizebox{0.72\textwidth}{!}{%
  \begin{tikzpicture}[
    ext/.style={rectangle, draw=UFSBlue, fill=UFSBlue!8, rounded corners,
      minimum width=3.2cm, minimum height=0.75cm, font=\small},
    int/.style={rectangle, draw=UFSYellow!80!black, fill=UFSYellow!25,
      minimum width=2.6cm, minimum height=0.7cm, font=\small},
    seta/.style={->, thick, UFSGray}
  ]
    \node[font=\Large\ttfamily] (expr) at (0,2.6) {a.nascimento.ano};

    \node[ext] (a)  at (-4.6,0.8) {\texttt{a}};
    \node[font=\scriptsize, align=center] at (-4.6,-0.1) {a variável\\\texttt{struct Aluno}};

    \node[ext] (n1) at (0,1.4)  {\texttt{nome[20]}};
    \node[ext] (n2) at (0,0.6)  {\texttt{matricula}};
    \node[ext, fill=UFSYellow!20] (n3) at (0,-0.2) {\texttt{nascimento}};
    \node[ext] (n4) at (0,-1.0) {\texttt{media}};
    \node[font=\scriptsize] at (0,-1.8) {campos de \texttt{struct Aluno}};

    \node[int] (d1) at (4.6,0.6)  {\texttt{dia}};
    \node[int] (d2) at (4.6,-0.2) {\texttt{mes}};
    \node[int, fill=UFSRed!15] (d3) at (4.6,-1.0) {\texttt{ano}};
    \node[font=\scriptsize] at (4.6,-1.8) {campos de \texttt{struct Data}};

    \draw[seta] (a) -- node[above=2pt, font=\scriptsize\ttfamily]{.nascimento} (n3);
    \draw[seta] (n3) -- node[above=2pt, font=\scriptsize\ttfamily]{.ano} (d3);
  \end{tikzpicture}}
  \end{center}
  \begin{block}{Regra}
    Cada ponto desce um nível. \texttt{a.nascimento} é um valor completo do tipo
    \texttt{struct Data}; \texttt{a.nascimento.ano} é o \texttt{int} dentro dele.
  \end{block}
\end{frame}

\begin{frame}[fragile]{Exemplo: usando o registro interno como valor}
\begin{lstlisting}[basicstyle=\ttfamily\tiny]
struct Data hoje = {17, 8, 2026};

struct Aluno a;
a.nascimento = hoje;              /* copia a struct Data inteira    */
a.nascimento.ano = 2005;          /* ajusta so um campo interno     */

/* funcao que serve a QUALQUER data, venha de onde vier */
int idade_em(struct Data nasc, struct Data ref) {
    int anos = ref.ano - nasc.ano;
    if (ref.mes < nasc.mes || (ref.mes == nasc.mes && ref.dia < nasc.dia))
        anos--;
    return anos;
}

printf("%d anos\n", idade_em(a.nascimento, hoje));   /* 21 anos */
\end{lstlisting}
  \begin{exampleblock}{Repare}
    \texttt{a.nascimento} entrou na função como argumento: o campo interno é
    um valor completo.
  \end{exampleblock}
\end{frame}

\begin{frame}[fragile]{Inicialização de registros aninhados}
\begin{lstlisting}[basicstyle=\ttfamily\scriptsize]
/* chaves internas espelham o aninhamento */
struct Aluno a = {"Ana Souza", 202600123, {7, 3, 2005}, 8.75};

/* com campos nomeados, fica auto-explicativo */
struct Aluno b = {
    .nome       = "Bia Rocha",
    .matricula  = 202600124,
    .nascimento = {.dia = 12, .mes = 11, .ano = 2004},
    .media      = 6.5
};
\end{lstlisting}
  \begin{itemize}
    \item As chaves internas são opcionais para o compilador, mas obrigatórias
      para quem vai ler o código depois.
    \item Em vetores: \texttt{turma[i].nascimento.ano} --- índice primeiro,
      depois os pontos.
  \end{itemize}
\end{frame}

\begin{frame}{Recapitulando: registros aninhados}
  \begin{block}{O que aprendemos}
    \begin{itemize}
      \item Um registro pode ser campo de outro, desde que declarado antes.
      \item Cada ponto desce um nível: \texttt{a.nascimento.ano}.
      \item O campo interno é um valor completo: pode ser atribuído de uma vez e
        passado para funções.
      \item Aninhar concentra um conceito (``data'') em um lugar só e permite
        reutilizar as funções que operam sobre ele.
    \end{itemize}
  \end{block}
\end{frame}

% =====================================================================
\section{Alinhamento e padding na memória}

\begin{frame}{De volta à pergunta: 44 ou 48?}
  \begin{columns}[T]
    \begin{column}{0.48\textwidth}
      \textbf{O que contamos:}\\[2mm]
      \small
      \begin{tabular}{lr}
        \toprule
        Campo & Bytes \\
        \midrule
        \texttt{char nome[20]}       & 20 \\
        \texttt{int matricula}       & 4 \\
        \texttt{struct Data nascimento} & 12 \\
        \texttt{double media}        & 8 \\
        \midrule
        \textbf{Soma}                & \textbf{44} \\
        \bottomrule
      \end{tabular}
    \end{column}
    \begin{column}{0.48\textwidth}
      \textbf{O que o compilador respondeu:}\\[2mm]
      \begin{center}
        \texttt{sizeof(struct Aluno) = 48}
      \end{center}
      \vspace{2mm}
      \small
      A diferença não é bug: o compilador \textbf{inseriu bytes} de propósito,
      para que cada campo comece em um endereço que o processador lê de uma vez.
    \end{column}
  \end{columns}
  \begin{block}{Vocabulário de hoje}
    \textbf{Alinhamento} --- em que endereços um tipo pode começar.
    \textbf{Offset} --- a que distância do início da struct o campo começa.
    \textbf{Padding} --- bytes vazios inseridos para o offset dar certo.
  \end{block}
\end{frame}

\begin{frame}{Por que o hardware exige alinhamento}
  \begin{center}
  \resizebox{0.78\textwidth}{!}{%
  \begin{tikzpicture}[
    byte/.style={rectangle, draw=UFSGray, minimum width=0.8cm,
      minimum height=0.7cm, font=\tiny},
    dado/.style={rectangle, draw=UFSBlue, fill=UFSBlue!25, minimum width=0.8cm,
      minimum height=0.7cm, font=\tiny},
    ruim/.style={rectangle, draw=UFSRed, fill=UFSRed!25, minimum width=0.8cm,
      minimum height=0.7cm, font=\tiny}
  ]
    % memoria alinhada
    \node[font=\small\bfseries, anchor=west] at (-2.6,1.9) {\texttt{int} alinhado (offset 4)};
    \foreach \i in {0,...,7} { \node[byte] at (\i*0.8,1.2) {}; }
    \foreach \i in {4,...,7} { \node[dado] at (\i*0.8,1.2) {}; }
    \foreach \x in {-0.4, 2.8, 6.0} {
      \draw[line width=1.4pt, UFSYellow!70!black] (\x,0.68) -- (\x,1.72);
    }
    \node[font=\scriptsize, anchor=west, text=UFSBlue] at (6.4,1.2)
      {\textbf{1} leitura de palavra};

    % memoria desalinhada
    \node[font=\small\bfseries, anchor=west] at (-2.6,0.1) {\texttt{int} desalinhado (offset 2)};
    \foreach \i in {0,...,7} { \node[byte] at (\i*0.8,-0.6) {}; }
    \foreach \i in {2,...,5} { \node[ruim] at (\i*0.8,-0.6) {}; }
    \foreach \x in {-0.4, 2.8, 6.0} {
      \draw[line width=1.4pt, UFSYellow!70!black] (\x,-1.12) -- (\x,-0.08);
    }
    \node[font=\scriptsize, anchor=west, text=UFSRed] at (6.4,-0.6)
      {\textbf{2} leituras + junção};

    \node[font=\scriptsize, anchor=west] at (-2.6,-1.6)
      {As barras douradas marcam as fronteiras de palavra da memória (4 bytes).};
  \end{tikzpicture}}
  \end{center}
  \begin{block}{A regra do hardware}
    Um tipo de $N$ bytes quer começar em um endereço múltiplo de $N$. Em x86 o
    acesso desalinhado é lento; em várias arquiteturas ARM e RISC ele
    simplesmente \textbf{falha}. O compilador prefere gastar bytes a arriscar isso.
  \end{block}
\end{frame}

\begin{frame}{As três regras que o compilador segue}
  \begin{block}{Ao montar o registro, na ordem em que você escreveu os campos}
    \begin{enumerate}
      \item Cada campo tem um \textbf{alinhamento} $a$ --- nos tipos básicos,
        igual ao tamanho (\texttt{char} 1, \texttt{short} 2, \texttt{int} 4,
        \texttt{double} 8).
      \item O campo é colocado no \textbf{próximo offset múltiplo de $a$};
        os bytes pulados viram \textbf{padding interno}.
      \item O alinhamento da struct é o \textbf{maior} entre os de seus campos,
        e o tamanho total sobe até o próximo múltiplo dele --- o
        \textbf{padding final}.
    \end{enumerate}
  \end{block}
  \begin{exampleblock}{Por que a regra 3 existe}
    Em \texttt{struct Aluno turma[40]}, cada elemento começa
    \texttt{sizeof(struct Aluno)} bytes depois do anterior. Sem múltiplo de 8,
    o campo \texttt{media} de \texttt{turma[1]} cairia desalinhado.
  \end{exampleblock}
\end{frame}

\begin{frame}{O mapa de bytes de \texttt{struct Aluno}}
  \begin{center}
  \resizebox{0.95\textwidth}{!}{%
  \begin{tikzpicture}[
    blk/.style={rectangle, draw=UFSBlue, fill=UFSBlue!12, minimum height=1.1cm,
      font=\scriptsize, align=center},
    pad/.style={rectangle, draw=UFSRed, fill=UFSRed!15, minimum height=1.1cm,
      font=\scriptsize, align=center, pattern=north east lines,
      pattern color=UFSRed!40},
    marca/.style={font=\scriptsize\ttfamily, text=UFSGray}
  ]
    % escala: 0.28cm por byte
    \def\s{0.28}
    % nome[20]: 0..19
    \node[blk, minimum width=20*\s cm, anchor=west] (n) at (0,0)
      {\texttt{nome[20]}\\\scriptsize 20 bytes};
    % matricula: 20..23
    \node[blk, minimum width=4*\s cm, anchor=west] (m) at (20*\s,0)
      {\texttt{matricula}\\\scriptsize 4};
    % nascimento: 24..35
    \node[blk, fill=UFSYellow!25, minimum width=12*\s cm, anchor=west] (d) at (24*\s,0)
      {\texttt{nascimento}\\\scriptsize 12 (3 $\times$ int)};
    % padding 36..39
    \node[pad, minimum width=4*\s cm, anchor=west] (p) at (36*\s,0) {};
    % media 40..47
    \node[blk, minimum width=8*\s cm, anchor=west] (v) at (40*\s,0)
      {\texttt{media}\\\scriptsize 8};

    % reguas de offset
    \foreach \o in {0,20,24,36,40,48} {
      \draw[UFSGray] (\o*\s,-0.62) -- (\o*\s,-0.95);
      \node[marca, anchor=north] at (\o*\s,-0.95) {\o};
    }
    \node[marca, anchor=north west] at (0,-1.5) {offset (bytes)};

    % anotacao do padding
    \node[font=\scriptsize, text=UFSRed, align=center, anchor=south]
      at (38*\s,0.75) {4 bytes de padding};
    \draw[->, UFSRed] (38*\s,0.72) -- (38*\s,0.58);
    \node[font=\scriptsize, text=UFSRed, align=left, anchor=west]
      at (49*\s,0) {\texttt{media} é \texttt{double}:\\precisa começar em\\múltiplo de 8 $\Rightarrow$ 40};
  \end{tikzpicture}}
  \end{center}
  \begin{block}{Conferindo}
    Offsets: \texttt{nome} 0, \texttt{matricula} 20, \texttt{nascimento} 24,
    \texttt{media} 40. Último byte usado: 47. Tamanho 48, múltiplo de 8 ---
    aqui não sobrou padding final, só os 4 bytes internos antes de \texttt{media}.
  \end{block}
\end{frame}

\begin{frame}[fragile]{Medindo em vez de adivinhar: \texttt{sizeof} e \texttt{offsetof}}
\begin{lstlisting}[basicstyle=\ttfamily\tiny]
#include <stdio.h>
#include <stddef.h>     /* offsetof */

int main(void) {
    printf("sizeof  = %zu\n", sizeof(struct Aluno));
    printf("nome       @ %zu\n", offsetof(struct Aluno, nome));
    printf("matricula  @ %zu\n", offsetof(struct Aluno, matricula));
    printf("nascimento @ %zu\n", offsetof(struct Aluno, nascimento));
    printf("media      @ %zu\n", offsetof(struct Aluno, media));
    return 0;
}
\end{lstlisting}
  \begin{exampleblock}{Saída (Linux/macOS 64 bits)}
    \texttt{sizeof = 48} \quad
    \texttt{nome @ 0} \quad \texttt{matricula @ 20} \quad
    \texttt{nascimento @ 24} \quad \texttt{media @ 40}
  \end{exampleblock}
  \begin{alertblock}{Não decore números}
    Os tamanhos dependem da arquitetura: aprenda o \emph{mecanismo} e meça.
  \end{alertblock}
\end{frame}

\begin{frame}[fragile]{A ordem dos campos muda o tamanho}
\begin{columns}[T]
\begin{column}{0.5\textwidth}
\begin{lstlisting}[basicstyle=\ttfamily\scriptsize]
struct Ruim {
    char c;
    int  i;
    char d;
};
/* sizeof = 12 */
\end{lstlisting}
\end{column}
\begin{column}{0.5\textwidth}
\begin{lstlisting}[basicstyle=\ttfamily\scriptsize]
struct Boa {
    int  i;
    char c;
    char d;
};
/* sizeof = 8 */
\end{lstlisting}
\end{column}
\end{columns}
  \begin{center}
  \resizebox{0.66\textwidth}{!}{%
  \begin{tikzpicture}[
    b/.style={rectangle, draw=UFSGray, minimum width=0.85cm, minimum height=0.65cm,
      font=\scriptsize\ttfamily},
    u/.style={b, draw=UFSBlue, fill=UFSBlue!18},
    p/.style={b, draw=UFSRed, fill=UFSRed!18}
  ]
    \node[font=\small\ttfamily, anchor=east] at (-0.6,0.7) {Ruim};
    \foreach \i/\t/\st in {0/c/u, 1/{}/p, 2/{}/p, 3/{}/p,
                           4/i/u, 5/i/u, 6/i/u, 7/i/u,
                           8/d/u, 9/{}/p, 10/{}/p, 11/{}/p}
      { \node[\st] at (\i*0.85,0.7) {\t}; }
    \node[font=\scriptsize, anchor=west, text=UFSRed] at (11.6,0.7) {6 bytes desperdiçados};

    \node[font=\small\ttfamily, anchor=east] at (-0.6,-0.4) {Boa};
    \foreach \i/\t/\st in {0/i/u, 1/i/u, 2/i/u, 3/i/u,
                           4/c/u, 5/d/u, 6/{}/p, 7/{}/p}
      { \node[\st] at (\i*0.85,-0.4) {\t}; }
    \node[font=\scriptsize, anchor=west, text=UFSBlue] at (11.6,-0.4) {2 bytes de padding final};
  \end{tikzpicture}}
  \end{center}
  \begin{block}{Regra prática}
    Declare os campos do \textbf{maior alinhamento para o menor}. Os mesmos
    dados ocuparam $33\%$ menos memória --- em um milhão de registros, isso
    são megabytes.
  \end{block}
\end{frame}

\begin{frame}{Quando isso importa de verdade}
  \begin{columns}[T]
    \begin{column}{0.48\textwidth}
      \textbf{Importa}
      \begin{itemize}
        \item Vetores grandes --- o padding é multiplicado pelo número de
          elementos.
        \item Sistemas embarcados, com RAM escassa.
        \item Desempenho: mais registros por linha de cache.
      \end{itemize}
    \end{column}
    \begin{column}{0.48\textwidth}
      \textbf{Não importa}
      \begin{itemize}
        \item Uma ou poucas variáveis do tipo.
        \item Quando reordenar prejudica a leitura sem ganho medido.
      \end{itemize}
    \end{column}
  \end{columns}
  \begin{alertblock}{Armadilha clássica}
    \texttt{fwrite(\&a, sizeof a, 1, f)} grava \emph{também} o padding, com
    lixo dentro --- outra máquina pode ler tudo errado.
  \end{alertblock}
\end{frame}

\begin{frame}{Recapitulando: alinhamento e padding}
  \begin{block}{O que aprendemos}
    \begin{itemize}
      \item \texttt{sizeof} de um registro $\geq$ soma dos campos: a diferença é
        \textbf{padding}.
      \item Cada campo começa em um offset múltiplo do seu alinhamento; o
        tamanho total é múltiplo do maior alinhamento da struct.
      \item O padding existe porque o processador lê a memória em blocos
        alinhados --- é preço de velocidade e de correção.
      \item Reordenar campos do maior para o menor reduz o desperdício, e
        \texttt{sizeof}/\texttt{offsetof} medem o resultado.
    \end{itemize}
  \end{block}
\end{frame}

% =====================================================================
\section{Tarefas}

\begin{frame}{Altere o programa}
  \begin{block}{Tarefa 1 --- prever antes de compilar}
    Acrescente \texttt{char turma;} a \texttt{struct Aluno}, logo depois de
    \texttt{matricula}. Anote sua previsão de \texttt{sizeof} \emph{antes} de
    compilar; depois meça. Explique a diferença (ou a ausência dela).
  \end{block}
  \begin{block}{Tarefa 2 --- reordenar}
    Reordene os campos de \texttt{struct Aluno} do maior alinhamento para o
    menor e meça \texttt{sizeof} antes e depois. Deu ganho? Depois repita com
    \texttt{struct \{ char flag; double valor; int id; char tipo; \}} e explique
    por que os dois casos se comportam de forma diferente.
  \end{block}
  \begin{block}{Tarefa 3 --- aninhar mais um nível}
    Crie \texttt{struct Endereco} (rua, número, CEP) e inclua-a em
    \texttt{struct Aluno}. Imprima \texttt{a.endereco.cep} e recalcule o
    tamanho total.
  \end{block}
\end{frame}

\begin{frame}{Desafio}
  \begin{exampleblock}{Cadastro da turma}
    Escreva um programa que:
    \begin{enumerate}
      \item declare \texttt{struct Aluno turma[5]} com data aninhada;
      \item preencha os cinco alunos no próprio código;
      \item imprima uma tabela com nome, matrícula, nascimento e média;
      \item informe o aluno de maior média e a média da turma;
      \item imprima \texttt{sizeof(struct Aluno)}, a soma dos campos e quantos
        bytes de padding existem por aluno e no vetor inteiro.
    \end{enumerate}
  \end{exampleblock}
  \begin{block}{Entrega}
    O tutorial traz o esqueleto, os testes e as perguntas de investigação.
  \end{block}
\end{frame}

\begin{frame}{Resumo da aula}
  \begin{center}
  \begin{tikzpicture}[
    item/.style={rectangle, draw=UFSBlue, fill=UFSBlue!10, rounded corners,
      minimum width=10.5cm, minimum height=0.75cm, font=\small}
  ]
    \node[item] (t1) at (0,0)    {\texttt{struct} agrupa campos de tipos diferentes em um valor único};
    \node[item] (t2) at (0,-1.1) {O ponto acessa campos; cada ponto desce um nível de aninhamento};
    \node[item] (t3) at (0,-2.2) {O compilador insere padding para alinhar cada campo};
    \node[item] (t4) at (0,-3.3) {\texttt{sizeof} e \texttt{offsetof} medem; a ordem dos campos decide o custo};
    \draw[->, thick, UFSBlue] (t1)--(t2);
    \draw[->, thick, UFSBlue] (t2)--(t3);
    \draw[->, thick, UFSBlue] (t3)--(t4);
  \end{tikzpicture}
  \end{center}
  \begin{block}{Próxima aula}
    Ponteiros: em vez de copiar 48 bytes, passar um endereço --- e o operador
    \texttt{->} para chegar aos campos a partir dele.
  \end{block}
\end{frame}

% =====================================================================
\section*{Referências}
\begin{frame}[allowframebreaks]{Referências}
  \scriptsize
  \nocite{*}
  \bibliographystyle{plain}
  \bibliography{../referencias}
\end{frame}

\end{document}
